\documentclass[output=paper]{langscibook}
\ChapterDOI{10.5281/zenodo.22880270}
\title{\textit{O sinivi suuna} – The indigenous canoe}
\author{Mark Tahi and Enoch Horai Magum}
\abstract{}
\begin{document}
\maketitle\label{chap:22}
\begin{figure}[H]
    \includegraphics[trim=0 0 0 0,clip,width=.85\linewidth]{figures/22_1.jpg}
    \includegraphics[trim=0 0 0 0,clip,width=.85\linewidth]{figures/22_2.jpg}
    \caption{\textit{Sinivi} \textit{suuna} -- Indigenous canoe}
    \label{fig:22_1-2}
\end{figure}
\begin{xltabular}{\textwidth}{lXX}
    (\ref{ex:text16-1})&\textit{O sinivi suuna}&The indigenous canoe\\
    (\ref{ex:text16-2})&\textit{O matapaku no sinivi suuna na kao vonaen:}&The method of making an indigenous canoe is as follows:\\
    (\ref{ex:text16-3})&\textit{Eara na hopo rara rasuu to tei nana o naono,}&We go into the jungle where the tree is,\\
    (\ref{ex:text16-4})&\textit{eara repaa sisira o naono havee toro busu voe beara tagava e.}&then we examine where the tree should fall down when we fell it.\\
    (\ref{ex:text16-5})&\textit{Eara repaa hanakahi vamataa nie.}&We clean the undergrowth completely for it.\\
    (\ref{ex:text16-6})&\textit{Be kuu gunaha, eara repaa ipinoo e, ei teo manin tea otei vai to pasi paku nana bono sinivi.}&When (the tree) has fallen down, we measure it according to the wish of the man who will make the canoe.\\
    (\ref{ex:text16-7})&\textit{O ipinoo he na paku raara bona bebeahu na kuriraara, evehe ei kou ei, eara toro siisira vamataa o naono,}&The measurement we take with the length of our arm, while at the same time we must thoroughly examine the tree,\\
    (\ref{ex:text16-8})&\textit{eara mene mosi vahoara o hum vai o vageegeno ge to tanove nana.}&so that we won't accidentally cut a crooked part or (a part) that is bulky.\\
    (\ref{ex:text16-9})&\textit{Be kavara o hum kiu bona, eara repaa vaobete vamataa e.}&When this part of work is finished, we lay it down properly.\\
\end{xltabular}
\begin{figure}[H]
    \includegraphics[trim=0 0 0 0,clip,width=.9\linewidth]{figures/22_3.png}
    \caption{\textit{O naono tagava} -- The felled tree}
    \label{fig:22_3}
\end{figure}
\begin{figure}[H]
    \includegraphics[trim=0 0 0 0,clip,width=.9\linewidth]{figures/22_4.jpg}
    \caption{\textit{Poana} -- Slice the upper part off}
    \label{fig:22_4}
\end{figure}
\begin{figure}[H]
    \raggedleft
    \includegraphics[trim=0 0 0 0,clip,width=.9\linewidth]{figures/22_5.jpg}
    \caption{\textit{Paku a bua govara} -- Make the two pointed ends}
    \label{fig:22_5}
\end{figure}
\begin{xltabular}{\textwidth}{lXX}
    (\ref{ex:text16-10})&\textit{Murinae eara repaa poana mae,}&After that we slice its upper part off,\\
    (\ref{ex:text16-11})&\textit{eara repaa paku pete a kiu vaa tea papaesi.}&then we do the next work of cutting the sides into shape.\\
    (\ref{ex:text16-12})&\textit{Beara papaes vakavara, eara repaa vahio vamataa vahabana e,}&When we have finished cutting the sides, we place it properly again,\\
    (\ref{ex:text16-13})&\textit{eara repaa taneo tea kepo.}&then we start carving.\\
    (\ref{ex:text16-14})&\textit{Kepo, are kepo, beara kepo ae mosi gunagunaha vaantee bata maara o ipinoo ge o manin tea too vaahito no sinivi.}&Carve, then we carve, while we are carving and cutting downwards according to the measurement or the wish of the owner of the canoe.\\
    (\ref{ex:text16-15})&\textit{Murina maa taba vaarii eara repaa vaakopo e.}&After these things, we turn it upside down.\\
    (\ref{ex:text16-16})&\textit{Eara repaa paku mana maa a bua govara vaa matana ae vaa itao.}&Then we eventually make the two pointed ends at the bow and at the stern.\\
    (\ref{ex:text16-17})&\textit{Beara paku vakavara a bua govara, o sinivi kou tavusu vai ei.}&When we have done the two pointed ends, the canoe is taking shape now.\\
    (\ref{ex:text16-18})&\textit{Be kavara, eara repaa paku pete a mosi na taeno sinivi, o hum to kahi hiohio ni nae o sinivi.}&When (this is) finished, we do the cutting of the bottom of the canoe (from the outside), the part on which the canoe will sit.\\
    (\ref{ex:text16-19})&\textit{Murina mosi bona, eara repaa paku pete a meha papaesi.}&After this cut, we do another cut along the sides of the canoe.\\
    (\ref{ex:text16-20})&\textit{A papaesi bona to kahi sue bata ka nom araara havee toro mosi kaakoo ara.}&This cut along the sides will tell us where we must cut.\\
    (\ref{ex:text16-21})&\textit{Teo vuri beara mosi suusuku bata maara amaa papaes vai, eara mee na taatagi bata maara a bua govara, vaa matana ae teo itao.}&When we do the next cuts along these cuts on the sides, we are also getting the two pointed ends at the bow and the stern into shape.\\
    (\ref{ex:text16-22})&\textit{Be kavara, eara repaa verete vahaa o sinivi are paa taneo tea koro a komano sinivi.}&When this is finished, we turn the canoe over again and then start cutting out the inside of the canoe.\\
    (\ref{ex:text16-23})&\textit{A maamihu kiu vai na paku vaagum ni raara.}&All these various kinds of work we do in a group.\\
    (\ref{ex:text16-24})&\textit{Beara koro vaantee a rahi voosu, o vaagum vai repaa pahin rahi voosu bari.}&When we have cut out enough for pulling (the canoe) home, this group will immediately pull it home.\\
    (\ref{ex:text16-25})&\textit{Beo sinivi suguna vaan, a too vaahito kou paku komana vai bona maa vagegeana kiu vai,}&When the canoe arrives in the village, the owner does the last pieces of work himself,\\
    (\ref{ex:text16-26})&\textit{a vaamono ae koro vagegeana bona komano sinivi.}&cutting more wood out and doing the last bits of cutting out the inside of the canoe.\\
    (\ref{ex:text16-27})&\textit{Tea taatagi bona suin no sinivi toro pakue bona vanaa}&To get the body of the canoe ready he must do it with the plane\\
    (\ref{ex:text16-28})&\textit{ae toro taatagi bona komano sinivi bona koro.}&and must get the inside of the canoe into shape with the gouging tool.\\
    (\ref{ex:text16-29})&\textit{Be kavara, ore paa vaobete vaavaha bono sinivi kahi bona roava ae a huan be ihuana koana bono bon pam.  }&When this is finished, he puts the canoe away from the sun and the rain, while waiting for the day of lashing (the outrigger to the canoe).\\
\end{xltabular}
\vspace{-.5cm}
\begin{figure}[H]
    \captionsetup[subfigure]{oneside,margin={0cm,0cm}}
    \subfloat[]{\includegraphics[trim=0 0 0 0,clip,width=.4\linewidth]{figures/22_6.png}\label{fig:22_6}}\hfil
    \subfloat[]{\includegraphics[trim=0 0 0 0,clip,width=.4\linewidth]{figures/22_8.png}\label{fig:22_8}}
    \caption{(a) \textit{Koro} -- Gouging tool, (b) \textit{Vanaa} -- Plane}
\end{figure}
\begin{figure}[H]
    \includegraphics[trim=0 0 0 0,clip,width=.4\linewidth]{figures/22_7.jpg}
    \caption{\textit{Koro} -- Cutting out the hull}
    \label{fig:22_7}
\end{figure}
    
\end{document}